% GENERATED FILE. Edit canonical lessons, then run npm run generate:books.
% canonical-source-hash: fnv1a64:056f6a832bb1c189
% canonical-lessons: RU-C531-besporyadok, RU-C531-khlam, RU-C531-kladovka, RU-C531-skladyvat, RU-C531-vybrasyvat

\chapter{Clutter and Clearing Out}
\label{ch:ru-clutter-and-clearing-out}
\hlchaptermodality{531}

\begin{chapteropening}
\textbf{By the end of this chapter:} \emph{I can talk about mess, junk and putting things away.}

Complete the last lesson of chapter 531: I can talk about mess, junk and putting things away.
\end{chapteropening}

\section[\texorpdfstring{\ru{беспорядок} (besporyádok)}{беспорядок (besporyádok)}]{\ru{беспорядок} (besporyádok) — a mess, disorder}

\label{lesson:RU-C531-besporyadok}



\begin{quote}
Before the new one: say the Russian for to hang out, then the Russian for to wash out.
\end{quote}

\subsection*{You'll want to know: \ru{беспорядок}}
\textbf{\ru{беспорядок}} (\emph{besporyádok}) — \textquotedblleft{}a mess, disorder\textquotedblright{}.

\begin{quote}
\ru{В} \ru{комнате} \ru{сына} \ru{вечный} \ru{беспорядок}.
\end{quote}

\begin{quote}
\emph{V kómnate sýna véchnyy besporyádok.}
\end{quote}

\begin{quote}
My son's room is permanently in a mess.
\end{quote}

Masculine; genitive \ru{беспорядка}.

\begin{cousinweb}
\emph{\ru{беспорядок}} is built on \emph{\ru{без}} 'without' and \emph{\ru{порядок}} 'order', from the root of \emph{\ru{ряд}} 'a row'; EtymDB lists \emph{\ru{беспорядок}} without a source.
\end{cousinweb}

\begin{grammarlens}[title={What you've built}]
One more word for telling, explaining and describing at B1.
\end{grammarlens}

\subsection*{Your turn}
\begin{itemize}
\raggedright
  \item \emph{Say it:} \emph{besporyádok}
  \item \emph{Say it:} \emph{besporyádok}, once more
  \item \emph{Say it:} a sentence of your own with \emph{besporyádok}
  \item \emph{Recall:} say the Russian for to hang out, then the Russian for to wash out, then say \emph{besporyádok} again
\end{itemize}

\begin{tcolorbox}[breakable,colback=teal!4,colframe=teal!35!black,title={Before you move on}]
What does \ru{беспорядок} mean? (\textquotedblleft{}A mess, disorder\textquotedblright{}.) Say it once more.
\end{tcolorbox}
\section[\texorpdfstring{\ru{хлам} (khlam)}{хлам (khlam)}]{\ru{хлам} (khlam) — junk}

\label{lesson:RU-C531-khlam}



\begin{quote}
Before the new one: say the Russian for to wash out, then the Russian for a mess.
\end{quote}

\subsection*{You'll want to know: \ru{хлам}}
\textbf{\ru{хлам}} (\emph{khlam}) — \textquotedblleft{}junk\textquotedblright{}.

\begin{quote}
\ru{На} \ru{балконе} \ru{много} \ru{старого} \ru{хлама}.
\end{quote}

\begin{quote}
\emph{Na balkóne mnógo stárogo khláma.}
\end{quote}

\begin{quote}
There's a lot of old junk on the balcony.
\end{quote}

Masculine, a collective noun.

\begin{cousinweb}
EtymDB has no entry for \emph{\ru{хлам}}, and the origin is not settled.
\end{cousinweb}

\begin{grammarlens}[title={What you've built}]
One more word for telling, explaining and describing at B1.
\end{grammarlens}

\subsection*{Your turn}
\begin{itemize}
\raggedright
  \item \emph{Say it:} \emph{khlam}
  \item \emph{Say it:} \emph{khlam}, once more
  \item \emph{Say it:} a sentence of your own with \emph{khlam}
  \item \emph{Recall:} say the Russian for to wash out, then the Russian for a mess, then say \emph{khlam} again
\end{itemize}

\begin{tcolorbox}[breakable,colback=teal!4,colframe=teal!35!black,title={Before you move on}]
What does \ru{хлам} mean? (\textquotedblleft{}Junk\textquotedblright{}.) Say it once more.
\end{tcolorbox}
\section[\texorpdfstring{\ru{кладовка} (kladóvka)}{кладовка (kladóvka)}]{\ru{кладовка} (kladóvka) — a storeroom}

\label{lesson:RU-C531-kladovka}



\begin{quote}
Before the new one: say the Russian for a mess, then the Russian for junk.
\end{quote}

\subsection*{You'll want to know: \ru{кладовка}}
\textbf{\ru{кладовка}} (\emph{kladóvka}) — \textquotedblleft{}a storeroom\textquotedblright{}.

\begin{quote}
\ru{Пылесос} \ru{стоит} \ru{в} \ru{кладовке}.
\end{quote}

\begin{quote}
\emph{Pylesós stoít v kladóvke.}
\end{quote}

\begin{quote}
The vacuum cleaner is in the storeroom.
\end{quote}

Feminine, colloquial for \emph{\ru{кладовая}}.

\begin{cousinweb}
\emph{\ru{кладовка}} is built on \emph{\ru{кладовая}} 'a storeroom', from the root of \emph{\ru{класть}} 'to put'; EtymDB has no entry.
\end{cousinweb}

\begin{grammarlens}[title={What you've built}]
One more word for telling, explaining and describing at B1.
\end{grammarlens}

\subsection*{Your turn}
\begin{itemize}
\raggedright
  \item \emph{Say it:} \emph{kladóvka}
  \item \emph{Say it:} \emph{kladóvka}, once more
  \item \emph{Say it:} a sentence of your own with \emph{kladóvka}
  \item \emph{Recall:} say the Russian for a mess, then the Russian for junk, then say \emph{kladóvka} again
\end{itemize}

\begin{tcolorbox}[breakable,colback=teal!4,colframe=teal!35!black,title={Before you move on}]
What does \ru{кладовка} mean? (\textquotedblleft{}A storeroom\textquotedblright{}.) Say it once more.
\end{tcolorbox}
\section[\texorpdfstring{\ru{складывать} (skládyvat)}{складывать (skládyvat)}]{\ru{складывать} (skládyvat) — to fold; to put together}

\label{lesson:RU-C531-skladyvat}



\begin{quote}
Before the new one: say the Russian for junk, then the Russian for a storeroom.
\end{quote}

\subsection*{You'll want to know: \ru{складывать}}
\textbf{\ru{складывать}} (\emph{skládyvat}) — \textquotedblleft{}to fold; to put together\textquotedblright{}.

\begin{quote}
\ru{Папа} \ru{аккуратно} \ru{складывает} \ru{рубашки}.
\end{quote}

\begin{quote}
\emph{Pápa akkurátno skládyvayet rubáshki.}
\end{quote}

\begin{quote}
Dad folds his shirts neatly.
\end{quote}

Imperfective; perfective \ru{сложить}. Also to add up numbers.

\begin{cousinweb}
\emph{\ru{складывать}} is built on \emph{\ru{с}-} 'together' and the root of \emph{\ru{класть}} 'to put'; EtymDB lists the verb without a source.
\end{cousinweb}

\begin{grammarlens}[title={What you've built}]
One more word for telling, explaining and describing at B1.
\end{grammarlens}

\subsection*{Your turn}
\begin{itemize}
\raggedright
  \item \emph{Say it:} \emph{skládyvat}
  \item \emph{Say it:} \emph{skládyvat}, once more
  \item \emph{Say it:} a sentence of your own with \emph{skládyvat}
  \item \emph{Recall:} say the Russian for junk, then the Russian for a storeroom, then say \emph{skládyvat} again
\end{itemize}

\begin{tcolorbox}[breakable,colback=teal!4,colframe=teal!35!black,title={Before you move on}]
What does \ru{складывать} mean? (\textquotedblleft{}To fold; to put together\textquotedblright{}.) Say it once more.
\end{tcolorbox}
\section[\texorpdfstring{\ru{выбрасывать} (vybrásyvat)}{выбрасывать (vybrásyvat)}]{\ru{выбрасывать} (vybrásyvat) — to throw away}

\label{lesson:RU-C531-vybrasyvat}



\begin{quote}
Before the new one: say the Russian for a storeroom, then the Russian for to fold.
\end{quote}

\subsection*{You'll want to know: \ru{выбрасывать}}
\textbf{\ru{выбрасывать}} (\emph{vybrásyvat}) — \textquotedblleft{}to throw away\textquotedblright{}.

\begin{quote}
\ru{Каждую} \ru{весну} \ru{мы} \ru{выбрасываем} \ru{старые} \ru{вещи}.
\end{quote}

\begin{quote}
\emph{Kázhduyu vesnú my vybrásyvayem stáryye véshchi.}
\end{quote}

\begin{quote}
Every spring we throw out old things.
\end{quote}

Imperfective; perfective \ru{выбросить} (\ru{выброшу}).

\begin{cousinweb}
\emph{\ru{выбрасывать}} is built on \emph{\ru{вы}-} 'out' and \emph{\ru{бросать}} 'to throw'; EtymDB lists the verb without a source.
\end{cousinweb}

\begin{grammarlens}[title={What you've built}]
One more word for telling, explaining and describing at B1.
\end{grammarlens}

\subsection*{Your turn}
\begin{itemize}
\raggedright
  \item \emph{Say it:} \emph{vybrásyvat}
  \item \emph{Say it:} \emph{vybrásyvat}, once more
  \item \emph{Say it:} a sentence of your own with \emph{vybrásyvat}
  \item \emph{Recall:} say the Russian for a storeroom, then the Russian for to fold, then say \emph{vybrásyvat} again
\end{itemize}

\begin{tcolorbox}[breakable,colback=teal!4,colframe=teal!35!black,title={Before you move on}]
What does \ru{выбрасывать} mean? (\textquotedblleft{}To throw away\textquotedblright{}.) Say it once more.
\end{tcolorbox}
